\documentclass[a4paper,12pt]{article}
\usepackage[T2A]{fontenc}
\usepackage[utf8]{inputenc}
\usepackage[russian]{babel}
\usepackage{amsmath,amssymb,mathtools}
\usepackage[scaled=0.92]{DejaVuSans}
\renewcommand{\familydefault}{\sfdefault}
\usepackage{icomma}
\usepackage[margin=21mm,top=20mm,bottom=22mm]{geometry}
\usepackage{microtype}
\usepackage{tikz}
\usetikzlibrary{calc,arrows.meta,positioning,shapes.geometric}
\usepackage{tabularx,booktabs,array}
\usepackage{enumitem}
\usepackage{hyperref}
\usepackage{fancyhdr}
\usepackage{setspace}
\usepackage{xcolor}
\definecolor{accent}{HTML}{087A7B}
\definecolor{darkaccent}{HTML}{124D55}
\definecolor{textgray}{HTML}{30383E}
\definecolor{lightaccent}{HTML}{DCEDEC}
\definecolor{muted}{HTML}{6B7C83}
\color{textgray}
\hypersetup{hidelinks,pdftitle={Тригонометрия в задачах прикладного характера},pdfauthor={Артём Александрович Лёвин}}
\setstretch{1.08}
\setlength{\parindent}{0pt}
\setlength{\parskip}{3pt}
\setlist[itemize]{leftmargin=*,itemsep=2pt,topsep=3pt}
\setlist[enumerate]{leftmargin=*,itemsep=5pt,topsep=4pt}
\pagestyle{fancy}
\fancyhf{}
\fancyhead[L]{\small\color{muted} Тригонометрия · задачи прикладного характера}
\fancyhead[R]{\small\color{muted} 12.10.2026}
\fancyfoot[L]{\footnotesize\color{muted} Артём Александрович Лёвин}
\fancyfoot[R]{\footnotesize\color{muted} \thepage}
\renewcommand{\headrulewidth}{0.3pt}
\renewcommand{\footrulewidth}{0pt}
\setlength{\headheight}{15pt}
\newcommand{\sectiontitle}[1]{\vspace{4pt}{\Large\bfseries\color{darkaccent} #1}\par\vspace{1pt}{\color{accent}\rule{\linewidth}{0.55pt}}\vspace{3pt}}
\newcommand{\subsectiontitle}[1]{{\large\bfseries\color{darkaccent}#1}\par}
\newcommand{\key}[1]{{\color{darkaccent}\bfseries #1}}
\newcommand{\degc}{^{\circ}}
\newcommand{\answerline}{\hfill\rule{23mm}{0.35pt}}
\begin{document}
\hypersetup{pageanchor=false}
\begin{titlepage}
\thispagestyle{empty}
\begin{center}
\vspace*{24mm}
{\color{accent}\large МАТЕМАТИКА · ЕГЭ ПРОФИЛЬНОГО УРОВНЯ}\par
\vspace{24mm}
{\Huge\bfseries\color{darkaccent} Чек-лист}\par
\vspace{6mm}
{\LARGE\bfseries Тригонометрия в задачах\par прикладного характера}\par
\vspace{14mm}
{\color{accent}\rule{74mm}{0.8pt}}\par
\vspace{12mm}
{\large Из формулы — к углу.\par Из неравенства — к интервалу времени.}\par
\vspace{19mm}
{\large 12 октября 2026 года}\par
\vfill
{\normalsize\bfseries Лёвин Артём Александрович}\par
{\normalsize эксклюзивно для Ярослава Гаврилова}\par
\vspace{20mm}
\begin{tikzpicture}[x=1cm,y=1cm]
\draw[accent!65,line width=1.15pt] (-7,0) .. controls (-4,1.15) and (-2,-1) .. (0,0) .. controls (2,1) and (4,-1.15) .. (7,0);
\end{tikzpicture}
\end{center}
\end{titlepage}
\hypersetup{pageanchor=true}

\sectiontitle{01. Перед решением: правильно прочитайте условие}
\key{Первый шаг — выписать данные и расшифровать обозначения.} Для каждого символа полезно указать его физический смысл, затем отметить величину, которую требуется найти. Даже знакомая формула приводит к ошибке, если перепутаны $k$ (порядок максимума), $t$ (время), $T$ (период) и $2\alpha$ (удвоенный угол).

\begin{tabularx}{\linewidth}{@{}>{\bfseries\raggedright\arraybackslash}p{38mm}>{\raggedright\arraybackslash}X@{}}
\toprule
Формулировка & Как понимать математически \\
\midrule
«Не превышает $a$» & $x\le a$. Для наименьшего угла в задаче о решётке нужно обосновать выбор максимально допустимого $d$.\\
«Не менее $a$» / «не ниже $a$» & $x\ge a$. Если нужно найти границу допустимых углов, сначала рассматривают равенство.\\
«Превышает $a$» & $x>a$. Неравенство строгое; при определении длительности отдельные концы интервала на результат не влияют.\\
«Из первой секунды» & Рассматривают только $0\le t\le1$; всё за пределами отрезка отбрасывают.\\
«В процентах» & $\displaystyle \frac{\Delta t}{1\,\text{с}}\cdot100\%$. В ответе указывают число процентов.\\
\bottomrule
\end{tabularx}

\subsectiontitle{Тригонометрический минимум для занятия}
\begin{center}
\renewcommand{\arraystretch}{1.3}
\begin{tabular}{@{}lccccc@{}}
\toprule
Угол & $0\degc$ & $30\degc$ & $45\degc$ & $60\degc$ & $90\degc$\\
\midrule
$\sin\alpha$ & $0$ & $\frac12$ & $\frac{\sqrt2}{2}$ & $\frac{\sqrt3}{2}$ & $1$\\
$\tg\alpha$ & $0$ & $\frac{\sqrt3}{3}$ & $1$ & $\sqrt3$ & ---\\
\bottomrule
\end{tabular}
\end{center}
\key{Градусы и радианы.} В формуле $\sin(\pi t)$ аргумент дан в радианах, а в $\sin(120\degc t-30\degc)$ — в градусах (при $t$ в секундах). Угловые величины должны быть согласованы. Отдельно проверьте, нужен ли в ответе $\alpha$, $2\alpha$, значение энергии, доля или процент.

\clearpage
\sectiontitle{02. Из физической формулы — к углу}
\subsectiontitle{Пример 1. Дифракционная решётка}
Длина волны $\lambda=400$ нм, порядок максимума $k=2$, период решётки $d\le1600$ нм. Найдите минимальный угол $\varphi$ наблюдения максимума.
\[
 d\sin\varphi=k\lambda,
 \qquad \sin\varphi=\frac{2\cdot400}{d}.
\]
При $0\degc\le\varphi\le90\degc$ минимальный угол достигается при \key{наибольшем} допустимом периоде $d=1600$ нм. Тогда
\[
 \sin\varphi=\frac{800}{1600}=\frac12,
 \qquad \boxed{\varphi=30\degc}.
\]
\key{Проверка применимости:} $d\ge k\lambda$, иначе условие $\sin\varphi\le1$ не выполняется и максимум данного порядка невозможен.

\subsectiontitle{Пример 2. Энергия и удвоенный угол}
Пусть $Q=mv^2\sin^2\alpha$, $m=2$, $v=10$, и требуется минимальное $2\alpha$, при котором $Q\ge50$. Для физического угла $0\degc\le\alpha\le90\degc$:
\[
 200\sin^2\alpha\ge50
 \quad\Longrightarrow\quad \sin^2\alpha\ge\frac14
 \quad\Longrightarrow\quad \alpha\ge30\degc.
\]
Поэтому $\boxed{(2\alpha)_{\min}=60\degc}$. Здесь используется $\sin\alpha\ge0$ на указанном интервале. Если нужен минимальный угол, равенство определяет его границу.

\subsectiontitle{Маленькая, но важная проверка}
Если $\tg\alpha=1$ и $0\degc<\alpha<90\degc$, то $\alpha=45\degc$. Однако при просьбе найти $2\alpha$ ответ был бы $90\degc$. Не подменяйте искомую величину промежуточным результатом.

\sectiontitle{03. Два последовательных вычисления}
\subsectiontitle{Пример 3. Скорость колебаний и кинетическая энергия}
Скорость груза и его кинетическая энергия задаются формулами
\[
 v(t)=v_0\sin\!\left(\frac{2\pi t}{T}\right),
 \qquad E=\frac{mv^2}{2}.
\]
Известны $m=0{,}08$ кг, $v_0=0{,}5$ м/с, $T=12$ с, $t=1$ с. Сначала вычисляем скорость ($t$ — момент наблюдения, $T$ — период):
\[
 v(1)=0{,}5\sin\frac{2\pi}{12}
 =0{,}5\sin\frac\pi6=0{,}25\text{ м/с}.
\]
И только затем энергию:
\[
 E=\frac{0{,}08\cdot(0{,}25)^2}{2}
 =\boxed{0{,}0025\text{ Дж}}.
\]
\clearpage
\sectiontitle{04. Из неравенства — к доле времени}
\subsectiontitle{Пример 4. Скорость не ниже порога}
На первой секунде движения $v(t)=5\sin(\pi t)$ см/с. Какую долю этой секунды скорость была не менее $2{,}5$ см/с?
\[
5\sin(\pi t)\ge2{,}5
\quad\Longleftrightarrow\quad\sin(\pi t)\ge\frac12.
\]
На первой секунде $0\le\pi t\le\pi$, поэтому по единичной окружности
\[
 \frac\pi6\le\pi t\le\frac{5\pi}{6}
 \quad\Longleftrightarrow\quad
 \frac16\le t\le\frac56.
\]
\begin{center}
\begin{minipage}[c]{.34\linewidth}
\centering
\begin{tikzpicture}[scale=1.45]
\draw[muted!70,->](-1.35,0)--(1.35,0);
\draw[muted!70,->](0,-1.25)--(0,1.28);
\draw[accent,line width=.95pt](0,0)circle(1);
\draw[accent!60!black,line width=3.2pt]({cos(30)},{sin(30)})arc(30:150:1);
\draw[muted,dashed](-1.15,.5)--(1.15,.5);
\node[font=\scriptsize,above right]at(.8,.58){$\frac\pi6$};
\node[font=\scriptsize,above left]at(-.8,.58){$\frac{5\pi}6$};
\node[font=\scriptsize,right]at(0,.5){$\frac12$};
\end{tikzpicture}
\end{minipage}%
\begin{minipage}[c]{.64\linewidth}
\centering
\begin{tikzpicture}[x=7cm,y=1cm]
\draw[-{Stealth[length=2mm]},muted](-.02,0)--(1.11,0)node[below]{$t$};
\draw[accent,line width=4pt]({1/6},0)--({5/6},0);
\foreach \x/\lab in {0/$0$,0.1667/$\frac16$,0.8333/$\frac56$,1/$1$}{\draw[darkaccent](\x,-.12)--(\x,.12)node[below=3pt,font=\scriptsize]{\lab};}
\node[font=\scriptsize,above,accent]at(.5,.14){$\Delta t=\frac23$ с};
\end{tikzpicture}
\end{minipage}
\end{center}
Длительность $\Delta t=\frac56-\frac16=\frac23$ с. Искомая доля:
\[
\frac{\Delta t}{1\,\text{с}}=\frac23\approx\boxed{0{,}67}
\qquad\text{(округление до сотых).}
\]

\clearpage
\thispagestyle{fancy}
\vspace*{7mm}
\begin{center}
{\LARGE\bfseries\color{darkaccent} Алгоритм: как решать прикладную задачу}\par
\vspace{4mm}
{\color{muted} Найдите тип вопроса и пройдите по соответствующей ветви.}
\end{center}
\vspace{6mm}
\tikzset{
 flowstart/.style={ellipse,draw=accent,line width=.85pt,fill=lightaccent!55,minimum width=40mm,minimum height=12mm,align=center,font=\small\bfseries},
 flowact/.style={rounded corners=2.5pt,draw=accent,line width=.8pt,minimum width=54mm,text width=49mm,minimum height=15mm,align=center,inner sep=5pt,font=\small},
 flowcheck/.style={diamond,aspect=2.05,draw=accent,line width=.8pt,fill=lightaccent!45,text width=44mm,align=center,inner sep=2pt,font=\small\bfseries},
 flowarrow/.style={-{Stealth[length=2.5mm]},line width=.9pt,draw=darkaccent}
}
\begin{center}
\begin{tikzpicture}
\node[flowstart] (a) at (0,0) {Прочитайте условие};
\node[flowact] (b) at (0,-2.05) {Выпишите обозначения, значения и искомую величину};
\node[flowcheck] (c) at (0,-4.55) {Нужно найти долю времени?};
\node[flowact] (l1) at (-3.75,-7.15) {Подставьте данные в формулу, выделите $\sin\alpha$, $\sin^2\alpha$ или $\tg\alpha$};
\node[flowact] (l2) at (-3.75,-9.85) {Найдите угол или физическую величину; проверьте область допустимых значений};
\node[flowact] (l3) at (-3.75,-12.55) {Сравните с вопросом: нужен $\alpha$, $2\alpha$, $E$ или другой результат?};
\node[flowact] (r1) at (3.75,-7.15) {Запишите пороговое условие как неравенство для $\sin\theta$};
\node[flowact] (r2) at (3.75,-9.85) {Найдите угловые интервалы по окружности};
\node[flowact] (r3) at (3.75,-12.55) {Выразите $t$; пересеките найденные интервалы с $[0;1]$};
\node[flowact] (r4) at (3.75,-15.25) {Сложите длительности; найдите долю $\Delta t/1$ и, если нужно, проценты};
\node[flowstart] (end) at (0,-18.5) {Запишите ответ};
\draw[flowarrow](a)--(b);
\draw[flowarrow](b)--(c);
\draw[flowarrow](c.west)-|node[pos=.22,above,font=\footnotesize]{нет}(l1.north);
\draw[flowarrow](c.east)-|node[pos=.22,above,font=\footnotesize]{да}(r1.north);
\draw[flowarrow](l1)--(l2);
\draw[flowarrow](l2)--(l3);
\draw[flowarrow](r1)--(r2);
\draw[flowarrow](r2)--(r3);
\draw[flowarrow](r3)--(r4);
\draw[flowarrow](l3.south)|-(end.west);
\draw[flowarrow](r4.south)|-(end.east);
\end{tikzpicture}
\end{center}
\clearpage

\sectiontitle{05. Если у синуса есть фазовый сдвиг}
\subsectiontitle{Пример 5. Лампочка и напряжение}
Датчик работает по формуле
\[
 u(t)=2\sin(120\degc t-30\degc).
\]
Лампочка горит при $u(t)\ge1$ В. Найдите, какую часть первой секунды она горит (в процентах). 
\[
 \sin(120\degc t-30\degc)\ge\frac12.
\]
На нужном участке изменения аргумента подходят углы от $30\degc$ до $150\degc$:
\[
 30\degc\le120\degc t-30\degc\le150\degc
 \ \Longleftrightarrow\ \boxed{0{,}5\le t\le1{,}5}.
\]
Но условие спрашивает \key{только про первую секунду}. Поэтому пересекаем найденный отрезок с $[0;1]$:
\[
 [0{,}5;1{,}5]\cap[0;1]=[0{,}5;1].
\]
Время горения $\Delta t=1-0{,}5=0{,}5$ с; ответ $\boxed{50\%}$.

\begin{center}
\begin{tikzpicture}[x=7.6cm,y=1cm]
\draw[-{Stealth[length=2mm]},muted](-.03,0)--(1.63,0)node[right]{$t$};
\draw[line width=5pt,accent!30](0,0)--(1,0);
\draw[line width=3pt,accent](.5,0)--(1.5,0);
\draw[line width=4pt,darkaccent](.5,0)--(1,0);
\foreach \x/\lab in {0/$0$,0.5/$0{,}5$,1/$1$,1.5/$1{,}5$}{\draw[darkaccent](\x,-.12)--(\x,.12)node[below=2pt,font=\scriptsize]{\lab};}
\node[above,anchor=south,font=\scriptsize,muted]at(.3,.16){первая секунда};
\node[above,anchor=south,font=\scriptsize,darkaccent]at(.9,.55){учитываем $0{,}5$ с};
\end{tikzpicture}
\end{center}

\subsectiontitle{Пример 6. Условие выполнено уже при $t=0$}
Для $u(t)=2\sin(150\degc t+45\degc)$ лампочка загорается при $u(t)\ge1$ В. На отрезке $0\le t\le1$ аргумент изменяется от $45\degc$ до $195\degc$. Поэтому допустима часть диапазона от $45\degc$ до $150\degc$:
\[
 150\degc t+45\degc\le150\degc
 \quad\Longrightarrow\quad t\le\frac{105}{150}=0{,}7.
\]
Лампочка горит с начала наблюдения до $t=0{,}7$ с, то есть $\boxed{70\%}$ первой секунды.

\subsectiontitle{Контроль ошибок}
\begin{itemize}
\item \key{Неправильный интервал:} ответ $[0{,}5;1{,}5]$ ещё не даёт искомую долю первой секунды. Обязательно пересекайте с $[0;1]$.
\item \key{Пропущенные обороты:} в общих задачах исследуйте все подходящие периоды синуса. В примерах занятия диапазон аргумента на $[0;1]$ короткий, поэтому хватает одного выделенного участка.
\item \key{Путаница с процентами:} $0{,}5$ доли $=50\%$, но в бланке записывают число согласно требуемому формату.
\end{itemize}


\clearpage
\sectiontitle{06. Тренировочный блок · Путь В}
{\color{muted}\small Два учебных дня: 13 и 14 октября 2026 года. На каждый день --- 5 математических заданий на изученные тригонометрические формулы и 3 прикладные задачи по занятию. Решения оформляйте самостоятельно. Ответы приведены в конце пособия.}\par

\vspace{3mm}
\subsectiontitle{13 октября 2026 года}
\vspace{2mm}
{\bfseries\color{accent} Часть I. Формулы из пройденных тем (5 заданий)}
\begin{enumerate}[label=\textbf{\color{accent}\arabic*.},itemsep=7pt]
\item Восстановите основное тождество $\sin^2\alpha+\cos^2\alpha=\ldots$ и найдите $\cos^2\alpha$, если $\sin\alpha=\frac35$. \answerline
\item Запишите формулу $\sin(\alpha+\beta)$ и вычислите $\sin75\degc$. \answerline
\item Восстановите формулу $\sin2\alpha$ и найдите её значение, если $\sin\alpha=\frac35$, $\cos\alpha=\frac45$. \answerline
\item Запишите $\cos2\alpha$ через $\sin\alpha$ и вычислите $\cos2\alpha$ при $\sin\alpha=\frac14$. \answerline
\item По формулам приведения выразите через функции угла $\alpha$: $\sin\bigl(\frac\pi2+\alpha\bigr)$ и $\cos(\pi+\alpha)$. \answerline
\end{enumerate}
\vspace{3mm}
{\bfseries\color{accent} Часть II. Прикладные задачи (3 задания)}
\begin{enumerate}[label=\textbf{\color{accent}\arabic*.},start=6,itemsep=10pt]
\item На дифракционную решётку падает свет с длиной волны $500$ нм. Период решётки не превышает $2000$ нм. Под каким минимальным углом (в градусах) возможен максимум второго порядка? Используйте $d\sin\varphi=k\lambda$ и $0\degc\le\varphi\le90\degc$. \answerline
\item Величина $Q=mv^2\sin^2\alpha$ должна быть не меньше $75$, если $m=3$, $v=10$, $0\degc\le\alpha\le90\degc$. Найдите минимальное значение \emph{удвоенного угла} $2\alpha$ (в градусах). \answerline
\item Скорость тела $v(t)=9\sin(\pi t)$ см/с. Какую долю \emph{первой секунды} она была не менее $4{,}5$ см/с? Округлите ответ до сотых. \answerline
\end{enumerate}

\clearpage
\sectiontitle{07. Путь В · второй день}
{\color{muted}\small Путь В: ещё 5 упражнений на ранее изученные математические формулы и 3 задачи прикладного характера. При ответе учитывайте, какая именно величина запрошена и в каких единицах.}\par
\vspace{4mm}
\subsectiontitle{14 октября 2026 года}
\vspace{2mm}
{\bfseries\color{accent} Часть I. Формулы из пройденных тем (5 заданий)}
\begin{enumerate}[label=\textbf{\color{accent}\arabic*.},start=9,itemsep=7pt]
\item Восстановите формулу $\sin(\alpha-\beta)$ и вычислите $\sin15\degc$. \answerline
\item Запишите формулу $\cos(\alpha+\beta)$ и вычислите $\cos105\degc$. \answerline
\item Вспомните формулу для $\tg\alpha$ через синус и косинус. Найдите $\tg\alpha$, если $\sin\alpha=-\frac35$, $\cos\alpha=-\frac45$. \answerline
\item Запишите формулы наименьших положительных периодов $\sin(kx)$ и $\tg(kx)$. Найдите периоды $\sin\frac{5x}{4}$ и $\tg\frac{2x}{3}$. \answerline
\item По формулам двойного угла упростите $1-2\sin^2x$ и $2\sin x\cos x$. \answerline
\end{enumerate}
\vspace{3mm}
{\bfseries\color{accent} Часть II. Прикладные задачи (3 задания)}
\begin{enumerate}[label=\textbf{\color{accent}\arabic*.},start=14,itemsep=10pt]
\item Для $Q=mv^2\sin^2\alpha$ известно: $m=2$, $v=10$, $0\degc\le\alpha\le90\degc$. Найдите наименьшее значение $2\alpha$ (в градусах), при котором $Q\ge150$. \answerline
\item Скорость груза $v(t)=v_0\sin(2\pi t/T)$; кинетическая энергия $E=mv^2/2$. Найдите $E$ (в Дж), если $m=0{,}4$ кг, $v_0=0{,}5$ м/с, $T=12$ с, $t=2$ с. \answerline
\item Напряжение датчика $u(t)=5\sin(150\degc t+15\degc)$ В. Лампочка горит при напряжении \emph{не ниже} $2{,}5$ В. Сколько процентов \emph{первой секунды} она горит? \answerline
\end{enumerate}

\clearpage
\sectiontitle{08. Ответы и самопроверка}
\vspace{-1mm}
\subsectiontitle{13 октября · задания 1--8}
\renewcommand{\arraystretch}{1.33}
\begin{tabularx}{\linewidth}{@{}>{\bfseries}p{12mm}p{23mm}X@{}}
\toprule
№ & Ответ & Ключевой контрольный шаг \\
\midrule
1 & $\frac{16}{25}$ & $\sin^2\alpha+\cos^2\alpha=1$.\\
2 & $\frac{\sqrt6+\sqrt2}{4}$ & $\sin(a+b)=\sin a\cos b+\cos a\sin b$.\\
3 & $\frac{24}{25}$ & $\sin2\alpha=2\sin\alpha\cos\alpha$.\\
4 & $\frac78$ & $\cos2\alpha=1-2\sin^2\alpha$.\\
5 & $\cos\alpha$; $-\cos\alpha$ & $\sin(\frac\pi2+\alpha)=\cos\alpha$, $\cos(\pi+\alpha)=-\cos\alpha$.\\
6 & $30$ & $d_{\max}=2000$, $\sin\varphi=\frac12$.\\
7 & $60$ & $\sin^2\alpha\ge\frac14$, затем $2\alpha$.\\
8 & $0{,}67$ & $t\in[1/6;5/6]$, доля $2/3$.\\
\bottomrule
\end{tabularx}
\vspace{5mm}
\subsectiontitle{14 октября · задания 9--16}
\begin{tabularx}{\linewidth}{@{}>{\bfseries}p{12mm}p{23mm}X@{}}
\toprule
№ & Ответ & Ключевой контрольный шаг \\
\midrule
9 & $\frac{\sqrt6-\sqrt2}{4}$ & $\sin(a-b)=\sin a\cos b-\cos a\sin b$.\\
10 & $\frac{\sqrt2-\sqrt6}{4}$ & $\cos(a+b)=\cos a\cos b-\sin a\sin b$.\\
11 & $\frac34$ & $\tg\alpha=\frac{\sin\alpha}{\cos\alpha}$.\\
12 & $\frac{8\pi}{5}$; $\frac{3\pi}{2}$ & $T_{\sin(kx)}=\frac{2\pi}{|k|}$, $T_{\tg(kx)}=\frac\pi{|k|}$.\\
13 & $\cos2x$; $\sin2x$ & $1-2\sin^2x=\cos2x$, $2\sin x\cos x=\sin2x$.\\
14 & $120$ & $\sin^2\alpha\ge\frac34$, $\alpha_{\min}=60\degc$.\\
15 & $0{,}0375$ & $v^2=\frac3{16}$, $E=0{,}4\cdot\frac3{32}$.\\
16 & $80$ & $t\in[0{,}1;0{,}9]$; доля $0{,}8$.\\
\bottomrule
\end{tabularx}
\vspace{5mm}
\subsectiontitle{Перед сдачей работы}
{\small
\begin{itemize}[itemsep=1pt,topsep=2pt]
\item В формульных упражнениях проверьте знаки, аргументы и коэффициенты тождеств.
\item В прикладных задачах определите, нужен ли угол, $2\alpha$, энергия, доля или процент.
\item Проверьте единицы измерения и аргумент синуса (градусы или радианы).
\item В задачах о первой секунде учитывайте лишь $0\le t\le1$.
\item Не теряйте квадрат в $\sin^2\alpha$; округляйте в конце решения.
\end{itemize}}
\vfill
{\color{accent}\rule{\linewidth}{0.5pt}}\par
{\small\color{muted} Лёвин Артём Александрович · эксклюзивно для Ярослава Гаврилова · 12.10.2026}
\end{document}
